\chapter{Compact targets and gravitational comparisons}
\label{mht:ch:compact-gravity}

A formal boundary chart determines local fields and their operations.
A compact theory must also glue those charts, and a gravitational
theory must identify their observables with fields defined by its own
action. The comparison needs a map of field complexes compatible
with products and collisions. It also needs states on which
amplitudes can be sewn.

The local trace construction of
Chapter~\ref{mht:ch:matrix-branes} supplies one component.
The compact examples below have additional global degrees of freedom.
Their Hall products and root characters can be constructed on explicit
coefficient spaces. The gravitational examples likewise admit explicit
connection fields, representations and prepared states. These
constructions let us ask which operations a comparison must preserve.

\section{Calabi--Yau traces and compact fields}

Let $\mathcal C$ be a small stable $\mathbb C$-linear category.
Properness means that each morphism complex is perfect over
$\mathbb C$. Smoothness means that its diagonal bimodule is perfect.
Assume $\mathcal C$ is smooth and proper. A $d$-Calabi--Yau trace pairs
$\operatorname{RHom}(E,F)$ with $\operatorname{RHom}(F,E)$
nondegenerately in degree $-d$. We require a negative-cyclic lift
when using this trace in a closed differential form on a moduli space.

At an object $E$, the tangent complex to the derived moduli of objects
is $\operatorname{RHom}(E,E)[1]$. Composition and trace give the pairing
\[
 \operatorname{RHom}(E,E)[1]^{\otimes2}
 \longrightarrow \operatorname{RHom}(E,E)[2]
 \longrightarrow \mathbb C[2-d].
\]
On a moduli stack to which the shifted symplectic construction applies,
Calabi--Yau duality makes this pairing nondegenerate, and the cyclic
lift supplies its closedness. The shift $2-d$ belongs to this moduli
problem. Hochschild cochains have their own composition operations.
Identifying them with operations on physical fields requires a map.

A concrete compact field model starts with a K3 surface $S$, an
elliptic curve $E$, and holomorphic forms $\sigma_S$ and $dz$.
Then $X=S\times E$ has volume form $\Omega_X=\sigma_S\wedge dz$.
Let $\mathfrak l$ be a finite-dimensional Lie algebra and let
$\mathfrak d=\mathfrak l\ltimes\mathfrak l^*$ be its cotangent Lie
algebra. Its pairing is the symmetric evaluation pairing, and its
bracket uses the coadjoint action on $\mathfrak l^*$.

The holomorphic cotangent BF fields are
\[
 A\in\Omega^{0,*}(X,\mathfrak l)[1],\qquad
 B\in\Omega^{0,*}(X,\mathfrak l^*)[1],
\]
with action
\begin{equation}\label{mht:eq:compact-bf-action}
 S(A,B)=\int_X\Omega_X\wedge
       \left\langle B,\bar\partial A+\tfrac12[A,A]\right\rangle.
\end{equation}
The integral pairs complementary Dolbeault degrees. Including the
field shifts, it is a BV pairing of degree $-1$.

\begin{proposition}\label{mht:prop:compact-bf}
The action \eqref{mht:eq:compact-bf-action} satisfies the classical
master equation. Choose a K\"ahler metric on $S$ and a harmonic
$(0,2)$-form $\eta$ with
$\int_S\sigma_S\wedge\eta=1$. Put
$K=\mathbb C[\eta]/(\eta^2)$, with $|\eta|=2$ and zero differential.
For every open $U\subset E$, inclusion of harmonic representatives gives
a trace-preserving dg Lie quasi-isomorphism
\[
 K\otimes\Omega^{0,*}(U)\otimes\mathfrak d
 \longrightarrow \Omega^{0,*}(S\times U)\otimes\mathfrak d.
\]
The smooth tensor product in the $S$ direction is completed.
The maps commute with restriction in $U$.
\end{proposition}

\begin{proof}
The differential of the cotangent theory is the cotangent lift of
$\bar\partial A+\tfrac12[A,A]$. Its square vanishes by
$\bar\partial^2=0$, the derivation rule and Jacobi. Invariance of the
evaluation pairing identifies this identity with the classical master
equation.

The map $i:K\to\Omega^{0,*}(S)$ is multiplicative because
$(0,4)$-forms on $S$ vanish. The Dolbeault cohomology of $S$ is
one-dimensional in degrees zero and two and vanishes in degree one.
Hodge projection $p$ and $h=\bar\partial^*G$ therefore satisfy
\[
 pi=1,\qquad \bar\partial h+h\bar\partial=1-ip,\qquad hi=0.
\]

Tensor these continuous operators with the identity on
$\Omega^{0,*}(U)\otimes\mathfrak d$, with the tensor differential.
They give the required contraction. Multiplicativity of $i$ proves
bracket compatibility, and the normalization of $\eta$ proves trace
compatibility. Restriction acts on the $U$ factor alone.
\end{proof}

The two K3 modes both survive this comparison. All higher
transferred products in the K3 factor vanish. The first product of
harmonic inputs lies in $i(K)$, and the next internal homotopy kills it.
This proves a classical field comparison. Quantized observables still
require compatible kernels, local extensions and observable
completions.

There is a useful finite-dimensional control on loop coefficients.
Suppose $\mathfrak l$ has a strictly positive grading and is truncated
above a finite height. Each cubic vertex has one $B$-leg. Since an
internal edge consumes one such leg, a connected graph has at most one
loop.

At one loop every $B$-leg is internal. The Lie coefficient is a
trace of adjoint operators with positive total height, hence is zero.
The dual trace vanishes for the same reason. This removes those Lie
coefficients. Analytic extension of the graph kernels and the full
quantum master equation remain separate requirements.

\section{Extension flags and their coefficient theory}

Let $\mathcal M_\gamma$ be a moduli stack of objects of charge
$\gamma$. The extension stack $\mathcal E_{\alpha,\beta}$ records
exact sequences $0\to E_\alpha\to E\to E_\beta\to0$ and has maps
\[
 \mathcal M_\alpha\times\mathcal M_\beta
 \xleftarrow{\ p\ }\mathcal E_{\alpha,\beta}
 \xrightarrow{\ q\ }\mathcal M_{\alpha+\beta}.
\]
Choose a coefficient theory in which $p^!$ and $q_*$ are defined on
the selected classes. This choice includes support conditions and,
for critical coefficients, orientation isomorphisms. A Hall product
then has the form $q_*p^!$, with its declared shifts.

The associativity comparison lives on the stack of two-step flags.
If base change identifies the two iterated pull--push constructions
with that same flag correspondence, their products agree. For three
or more steps these identifications must be compatible under
forgetting intermediate terms. Thus orientation and base change are
part of the multiplication, rather than consequences of a numerical
count of extensions.

An action on a field theory requires one further correspondence map.
Its target is a state complex defined from the fields, independently
of the Hall action. For a compact chiral state space, this map must
retain the normal oscillators at diagonals and commute with joining
cycles. Local coordinate formulas must also agree through higher
changes of coordinates. Defining the state space to be the image of
the Hall map would omit this comparison problem.

\section{An exchange ratio does not determine the integral algebra}

The following polynomial example isolates the coefficient issue.
Put
\[
 \begin{gathered}
 R=\mathbb Q[h_1,h_2],\qquad h_3=-h_1-h_2,\\
 a=-(h_1^2+h_1h_2+h_2^2),\qquad b=-h_1h_2(h_1+h_2),\\
 N(v)=v^3+av+b.
 \end{gathered}
\]
The ring $R$ is a domain and $b\ne0$.

For $j=0,1$, take independent copies of
$\mathcal A^{(j)}=\bigoplus_{n\geq0}R[x_1,\ldots,x_n]^{S_n}$.
Charge is $n$. All cohomological degrees are zero and all parities are
even. The sum has finite charge support. No parameter or root
difference is inverted.

Define
\[
 K_0(v)=N(v)/v,\qquad K_1(v)=vN(v)=v^2K_0(v).
\]
and, for inputs of charges $d,e$, set
\begin{equation}\label{mht:eq:compact-shuffle}
 (F*_jG)(X)=
 \sum_{\substack{I\sqcup J=[d+e]\\|I|=d}}
 F(X_I)G(X_J)\prod_{i\in I,\,k\in J}K_j(x_k-x_i).
\end{equation}
There is no factorial divisor. The charge-zero constant is the unit.
The constant polynomial in charge one is a different element.

\begin{proposition}\label{mht:prop:compact-shuffle-lattices}
Both products are polynomial, unital and associative. They have the
same exchange ratio
\[
 r(v)=\frac{K_j(-v)}{K_j(v)}
     =\frac{v^3+av-b}{v^3+av+b},
\]
but no charge-preserving $R$-linear algebra isomorphism identifies
their full polynomial carriers. The conclusion also holds for maps
semilinear under one automorphism of $R$.
\end{proposition}

\begin{proof}
For $K_0$, pair summands across a possible pole $x_k=x_i$ by
exchanging the two labels. The remaining factors agree on that
diagonal and the two residues are opposite. All denominators have
simple root-difference factors, so cancellation proves polynomiality.
For $K_1$ polynomiality is immediate.

Both bracketings of three inputs
sum over the same ordered three-block partitions, with one kernel
factor per cross-block pair. This proves associativity in all charges.
Empty blocks prove the unit equations. The even factor $v^2$ cancels
from the exchange ratio.

In charge two write $c=(x+y)/2$, $v=y-x$, $t=v^2$ and $B=R[c,t]$.
Then $R[x,y]=B\oplus vB$. The linearized multiplication maps obey
\[
 \mu_0(A+vC)=2((t+a)A+bC),\qquad
 \mu_1(A+vC)=2t((t+a)A+bC).
\]
Their images are $J=(t+a,b)B$ and $tJ$. Hence
\[
 \operatorname{coker}\mu_0\simeq R[c]/(b),\qquad
 \operatorname{coker}\mu_1=B/tJ.
\]

The first cokernel has scalar annihilator $(b)$. The second surjects
onto the faithful $R$-module $R[c]$ by evaluation at $t=0$, so its
annihilator is zero. An algebra isomorphism preserving charge would
identify these cokernels. A semilinear isomorphism would take the
nonzero annihilator $b$ to its nonzero image under the scalar
automorphism. Both possibilities contradict faithfulness.
\end{proof}

The obstruction concerns the full polynomial realizations. A positive
comparison exists on their generated subalgebras. Put
\[
 D_n(X)=\prod_{i<k}(x_k-x_i)^2,\qquad D_0=D_1=1.
\]
The maps $\Phi_n(F)=D_nF$ are injective and multiplicative because
\[
 D_{d+e}(X)=D_d(X_I)D_e(X_J)
                   \prod_{i\in I,\,k\in J}(x_k-x_i)^2.
\]

They fix charge one. Their restrictions therefore identify the two
subalgebras generated by all charge-one polynomials. Inverting $D_n$
also makes them isomorphisms on the localized full carriers. Neither
statement gives an invertible map on the original full lattice:
the constant polynomial in charge $n\geq2$ is not divisible by $D_n$.

\section{Cooperation and the tensor of its target}

The same algebras expose a separate obstruction to an ordinary
bialgebra. Let $e_m=x^m$ in charge one and let $\varepsilon$ be the
projection to charge zero. An ordinary coproduct here means an
$R$-linear, charge-preserving, unital algebra map
\[
 \Delta:\mathcal A^{(j)}\longrightarrow
                \mathcal A^{(j)}\otimes_R\mathcal A^{(j)}
\]
with the usual tensor multiplication and both counit equations for
$\varepsilon$. The same definition can be made on the subalgebra
generated by the $e_m$.

\begin{proposition}\label{mht:prop:compact-ordinary-cooperation}
Neither algebra admits this coproduct, even on its current-generated
subalgebra. Coassociativity is not needed for the obstruction.
\end{proposition}

\begin{proof}
The charge-zero counits force every $e_m$ to be primitive. Direct
substitution in \eqref{mht:eq:compact-shuffle} gives
\begin{equation}\label{mht:eq:compact-current-relation}
 [e_0,e_3]-3[e_1,e_2]+a[e_0,e_1]=b(e_0*_je_0).
\end{equation}
Indeed, for polynomials $f,g$ the commutator is
\[
 2b\,v^{2j}\frac{f(x)g(y)-g(x)f(y)}{v},\qquad v=y-x.
\]
The numerator for the first two terms of
\eqref{mht:eq:compact-current-relation}, divided by $v$, is $v^2$.
The remaining term contributes $a$, and
$e_0*_je_0=2v^{2j}(v^2+a)$.

In an ordinary tensor algebra a commutator of even primitive elements
is primitive. The charge-$(1,1)$ part of the coproduct of the left
side is therefore zero. That of the right side is
$2b\,e_0\otimes e_0$, which is nonzero over $R$.
This contradicts multiplicativity.
\end{proof}

This proof uses $b\ne0$. It makes no assertion on the fiber $b=0$.
The charge-$(1,1)$ module is free, so merely writing a derived tensor
over $R$ does not remove the displayed defect.

There is again a positive construction after changing the recipient.
For two disjoint root lists put
\[
 Q_j(X,Y)=\prod_{x\in X,\,y\in Y}K_j(y-x),\qquad
 \Delta_{d,e}^{\mathrm{rat}}F(X,Y)=\frac{F(X,Y)}{Q_j(X,Y)}.
\]
Work on the ordered root cover with all pairwise differences and the
zeros of $N(\pm v)$ inverted, and then take symmetry within each
list. Let $A$ have bidegree $(d,e)$ and $B$ have bidegree $(f,g)$.
Define their product by summing
\begin{equation}\label{mht:eq:compact-rational-tensor}
 \begin{split}
 &A(X_1,Y_1)B(X_2,Y_2)
 Q_j(X_1,X_2)Q_j(Y_1,Y_2)\\
 &\hspace{15mm}\cdot
   \prod_{y\in Y_1,\,x\in X_2}r(y-x)
 \end{split}
\end{equation}
over partitions of $X$ into sizes $d,f$ and of $Y$ into sizes $e,g$.
The last factor changes the interchange of the two middle blocks.

\begin{proposition}\label{mht:prop:compact-rational-cooperation}
Formula \eqref{mht:eq:compact-rational-tensor} is associative and
unital. With this target multiplication, the sum of
$\Delta_{d,e}^{\mathrm{rat}}$ is multiplicative and coassociative,
with counits given by empty lists.
\end{proposition}

\begin{proof}
For three factors, either bracketing uses the same partitions of
$X$ and $Y$. It contributes one $K_j$ for every ordered pair within
each list. It also contributes one $r$ whenever a root in an earlier
$Y$ block crosses a root in a later $X$ block.
These scalar factors commute, so the
terms agree.

To compare a split shuffle product with the product of two splits,
partition each input between $X$ and $Y$. Cross factors with the
order $X$ before $Y$ cancel against the denominator $Q_j(X,Y)$.
The opposite order leaves $K_j(x-y)/K_j(y-x)=r(y-x)$, exactly the
factor in \eqref{mht:eq:compact-rational-tensor}. This proves
multiplicativity term by term. Three-way splitting divides by
$Q_j(X,Y)Q_j(X,Z)Q_j(Y,Z)$ in either order. Empty lists give the
counits.
\end{proof}

The elementary relation now reads
\[
 (1-r(v))v(v^2+a)=b(1+r(v)),
\]
so the changed interchange cancels the earlier defect. But the split
of the charge-two constant is $v/N(v)$ for $j=0$ and $1/(vN(v))$
for $j=1$. These are not polynomial coefficients. A comparison to
supported collision operations must restore the information on the
removed divisors, including any derived excess classes.

\section{Root states and scalar automorphic products}

Let $V=\bigoplus_{\alpha>0}V_\alpha$ be a graded superspace with
finite-dimensional pieces. Choose a positive height so that every
height truncation contains finitely many basis vectors. The completed
Fock space is
\[
 \mathcal H_V=\widehat{\operatorname{Sym}}(V_{\bar0})
                  \widehat\otimes\widehat\Lambda(V_{\bar1}).
\]

Completion is by height. A root variable $x^\alpha$ records the
additive degree, and the parity operator defines the supertrace.
For finitely many roots, or in an absolute-convergence chamber,
\begin{equation}\label{mht:eq:compact-fock-character}
 \operatorname{str}_{\mathcal H_V}(x^{\mathrm{degree}})
       =\prod_{\alpha>0}(1-x^\alpha)^{-\operatorname{sdim}V_\alpha}.
\end{equation}
An even oscillator contributes a geometric series, while an odd one
contributes $1-x^\alpha$. Multiplication proves the formula at each
finite height. Coefficientwise stabilization, or absolute convergence,
gives the stated limit.

Applied to specified Borcherds root spaces, this is the Fock product
associated with the denominator identity. The Weyl monomial, the
coordinate pullback and any doubling of root multiplicities must
still be included in comparing automorphic normalizations. In the
Igusa example, the weight-five square-root branch and the weight-ten
cusp form are sections of different modular lines. The reciprocal
partition product is a further scalar convention.

The supertrace loses information even before geometry enters.
One even and one odd vector of the same degree contribute $1$ to
\eqref{mht:eq:compact-fock-character}. Their Fock space is nonetheless
nontrivial. Nor does a graded trace determine a bracket on $V$.
Consequently a product identity cannot construct the brane objects,
their extension operations or a map to an independently defined
compact state space.

\begin{problem}\label{mht:prob:compact-geometric-comparison}
For a compact Calabi--Yau target, construct reduced critical
coefficients with their relative unit, cyclic pairing and orientation
on all extension flags. Construct their action on the independently
defined chiral state complexes, preserving charge, normal oscillators,
diagonal joining and higher coordinate gluing. Compare the resulting
trace line and its normalized section with the specified automorphic
line. On a formal normal chart, require compatibility with the native
trace coordinate and supported operations.
\end{problem}

This problem retains the positive classical compact field reduction,
the Hall products on their declared carriers and the scalar product
formula. Its missing maps concern the geometric field and state
recipients, rather than the existence of those algebraic constructions.

\section{Gravitational fields and the invertible coframe}

The gravitational comparison begins with its own bulk fields.
Let $M$ be an oriented three-manifold and fix $\ell>0$. Take the
Lorentz form $\eta_{ab}$ and its oriented volume tensor
$\epsilon_{ab}{}^c$. The Cartan Lie algebra has generators $J_a,P_a$
and brackets
\[
 [J_a,J_b]=\epsilon_{ab}{}^cJ_c,\qquad
 [J_a,P_b]=\epsilon_{ab}{}^cP_c,\qquad
 [P_a,P_b]=\ell^{-2}\epsilon_{ab}{}^cJ_c.
\]
The invariant form pairs $J_a$ with $P_b$ by $\eta_{ab}$ and vanishes
on the two diagonal summands. A connection is
$\mathcal A=\omega^aJ_a+e^aP_a$.

Write
\[
 R^a=d\omega^a+\tfrac12\epsilon^a{}_{bc}\omega^b\wedge\omega^c,
 \qquad
 T^a=de^a+\epsilon^a{}_{bc}\omega^b\wedge e^c.
\]
Its curvature is
\begin{equation}\label{mht:eq:cartan-curvature}
 F_{\mathcal A}
 =\left(R^a+\frac1{2\ell^2}
                  \epsilon^a{}_{bc}e^b\wedge e^c\right)J_a+T^aP_a.
\end{equation}

\begin{proposition}\label{mht:prop:cartan-invertible}
On the locus where $e$ is an invertible coframe, flatness of
$\mathcal A$ is equivalent to torsion-free Einstein geometry with
cosmological constant $\Lambda=-\ell^{-2}$ and metric
$g_{\mu\nu}=\eta_{ab}e^a_\mu e^b_\nu$.
\end{proposition}

\begin{proof}
The $P$ component of \eqref{mht:eq:cartan-curvature} gives $T=0$.
Invertibility identifies $\omega$ with the metric spin connection.
The $J$ component then says that the metric has constant sectional
curvature $-\ell^{-2}$. In three dimensions the Riemann tensor is
determined by the Ricci tensor, so this is the Einstein equation.
Conversely a torsion-free Einstein coframe has this curvature, and
both components of $F_{\mathcal A}$ vanish.
\end{proof}

Degenerate coframes remain connections but do not give this metric
interpretation. The same restriction matters in the linear theory:
at a flat background, a variation $a=\delta\mathcal A$ obeys
$d_{\mathcal A}a=0$, with gauge variation
$a\mapsto a+d_{\mathcal A}\epsilon$. Flatness gives
$d_{\mathcal A}^2=0$. Splitting $a$ into $\delta\omega$ and $\delta e$
is the linearized Cartan complex. Passing to metric variables uses
the invertible coframe and quotient by infinitesimal Lorentz rotations.
Boundary conditions and the induced pairing must be carried through
this passage before it compares Hilbert spaces.

The usual Chern--Simons normalization sets $k=\ell/(4G)$, and its
classical Brown--Henneaux relation is $c=6k$. These equations specify
a classical normalization. A quantum coupling comparison must
intertwine the renormalized current and ghost operations, including
their allowed finite counterterms. A chosen scalar reparametrization
is not that map.

\section{A Weyl representation and its modular transform}

Fix $b>0$. On $L^2(\mathbb R,dP)$, the self-adjoint Weyl pair has
Schwartz-space formulas
\[
 \widehat\ell_A=2\pi bP,\qquad
 \widehat\ell_B=-\frac{ib}{2}\frac{d}{dP},\qquad
 [\widehat\ell_A,\widehat\ell_B]=i\pi b^2.
\]
Reflection acts by $P\mapsto-P$. The individual length operators
exchange parity, while their squares preserve the even subspace.
Identify that subspace with
$\mathcal H=L^2(\mathbb R_{\geq0},dP)$ by
$f\mapsto 2^{-1/2}f(|P|)$.

The full-line Fourier transform with kernel
$\sqrt2\,e^{4\pi iPP'}$ is unitary and squares to reflection.
Its restriction is therefore the unitary involution
\begin{equation}\label{mht:eq:gravity-cosine}
 (\mathcal Ff)(P)=2\sqrt2\int_0^\infty
                  \cos(4\pi PP')f(P')\,dP'.
\end{equation}
For $t>0$ the Gaussian integral gives
\begin{equation}\label{mht:eq:gravity-gaussian}
 \mathcal F(e^{-tP^2})
   =\sqrt{\frac{2\pi}{t}}\,e^{-4\pi^2P^2/t}.
\end{equation}
These are representation-theoretic constructions on the specified
Lebesgue space. A physical zero-mode quantization must identify its
measure and observables with this representation.

Put $Q=b+b^{-1}$ and $c=1+6Q^2$, so $c\geq25$.
For $q=e^{2\pi i\tau}$ with $\operatorname{Im}\tau>0$, let
$\eta(\tau)=q^{1/24}\prod_{n\geq1}(1-q^n)$ be the Dedekind function.
The Virasoro character functions in this normalization are
\[
 \chi_P(\tau)=\frac{q^{P^2}}{\eta(\tau)}.
\]

The vacuum character is the difference at the two imaginary momenta
$iQ/2$ and $i(b^{-1}-b)/2$:
\[
 \chi_{\mathrm{vac}}(\tau)
 =q^{-c/24}\prod_{n\geq2}(1-q^n)^{-1}.
\]
The difference removes the level-one descendant. Transforming the
two Gaussians yields
\[
 \chi_{\mathrm{vac}}(-1/\tau)
       =\int_0^\infty\rho(P)\chi_P(\tau)\,dP,
\]
where
\begin{equation}\label{mht:eq:gravity-cap-density}
 \begin{split}
 \rho(P)
 &=2\sqrt2\bigl(\cosh(2\pi QP)
                  -\cosh(2\pi(b^{-1}-b)P)\bigr)\\
 &=4\sqrt2\,\sinh(2\pi bP)\sinh(2\pi P/b).
 \end{split}
\end{equation}

This density is positive for $P>0$ and vanishes quadratically at zero.
Its exponential growth prevents it from being a vector of
$\mathcal H$. It does define a continuous functional on each
Gaussian-weighted space
$L^2(\mathbb R_{\geq0},e^{P^2/n}dP)$, by Cauchy--Schwarz.
Their union, with its final locally convex topology, is a test domain
for the uncollared linear cap.

\section{Positive collars and finite sewing}

Let $H_0$ be multiplication by $P^2$ on $\mathcal H$, and write
$C_t=e^{-tH_0}$ for $t\geq0$. This is a bounded strongly continuous
semigroup. Positive collars give actual vectors
\[
 v_{0,t}(P)=e^{-tP^2},\qquad
 v_{1,t}(P)=\rho(P)e^{-tP^2},\qquad t>0.
\]
The Gaussian dominates the exponential growth of $\rho$, so both
belong to $\mathcal H$. A preparation is the map
$\mathbb C\to\mathcal H$ taking $1$ to one of these vectors.
Its reflected cap is the Hilbert adjoint.

Use an inner product linear in its second argument and define
\[
 K_{ij}(s,t)=|v_{i,s}\rangle\langle v_{j,t}|,\qquad
 G_{ij}(u)=\int_0^\infty\rho(P)^{i+j}e^{-uP^2}\,dP,
 \quad i,j\in\{0,1\}.
\]
The first expression is a rank-one operator. The second records the
overlap at the total collar width.

\begin{proposition}\label{mht:prop:gravity-prepared-sewing}
For positive widths,
\[
 \begin{split}
 K_{ij}(s,t)K_{k l}(u,v)
     &=G_{jk}(t+u)K_{i l}(s,v),\\
 K_{ij}(s,t)^*&=K_{ji}(t,s),\\
 \operatorname{Tr}K_{ij}(s,t)&=G_{ji}(s+t).
 \end{split}
\]
Every finite product of these kernels is trace class. More generally,
a bounded operator word with a trace-class factor has a finite cyclic
trace. For a Hilbert--Schmidt operator $B$,
$\operatorname{Tr}(B^*B)=\|B\|_{\mathrm{HS}}^2\geq0$.
\end{proposition}

\begin{proof}
Composition of rank-one operators contracts the two middle vectors,
giving the first formula. The adjoint and trace formulas follow from
the same inner product. The trace norm of $|u\rangle\langle v|$ is
$\|u\|\|v\|$. Thus sums with
$\sum_n\|u_n\|\|v_n\|<\infty$ converge in trace norm.
The rank-one trace identity extends to this completion and gives
$\operatorname{Tr}(AB)=\operatorname{Tr}(BA)$ for bounded $A$ and
trace-class $B$. The last assertion is the definition of the
Hilbert--Schmidt norm.
\end{proof}

All three overlap functions have closed Gaussian expressions. Put
$A=2\pi Q$, $B=2\pi(b^{-1}-b)$ and
$J_u(k)=\sqrt\pi\,e^{k^2/(4u)}/(2\sqrt u)$. Then
\[
 \begin{aligned}
 G_{00}(u)&=J_u(0),\\
 G_{01}(u)&=2\sqrt2\bigl(J_u(A)-J_u(B)\bigr),\\
 G_{11}(u)&=4J_u(2A)+4J_u(2B)+8J_u(0)\\
          &\hspace{8mm}-8J_u(A+B)-8J_u(A-B).
 \end{aligned}
\]
The linear cap uses one factor of $\rho$. Reflected sewing of two
caps uses $\rho^2$. Replacing $\rho$ by its square root would change
the character transform \eqref{mht:eq:gravity-cap-density}.

Unitary conjugation and Hilbert tensor products preserve these
trace-class sewing identities. Remarking the torus conjugates the
cylinder as well as its states. In particular $\mathcal F$ does not
commute with $C_t$: applying \eqref{mht:eq:gravity-gaussian} to
$C_s e^{-tP^2}$ gives a different width from
$C_s\mathcal F e^{-tP^2}$. Thus unitary transport is not invariance
of a fixed cylinder operator.

\section{The trace obstruction for the continuous cylinder}

The finite amplitudes above use prepared vectors. Positive time alone
does not make $C_t$ trace class. On the infinite-dimensional subspace
$L^2([1,2])$, one has $C_t\geq e^{-4t}$. The sum of its diagonal
matrix coefficients over any $N$ orthonormal vectors in this subspace
is at least $Ne^{-4t}$. In particular $C_t$ is not compact.

The obstruction persists for a finite trace with no continuity or
positivity assumption, provided it is linear and cyclic on a full
operator ideal.

\begin{theorem}\label{mht:thm:gravity-no-finite-trace}
Let $t>0$ and let $I\subset B(\mathcal H)$ be a two-sided ideal
containing $C_t$. No complex-linear map $\tau:I\to\mathbb C$
satisfies both of the following conditions:
\[
 \begin{aligned}
 \tau(AB)&=\tau(BA)&& (A\in I,\ B\in B(\mathcal H)),\\
 \tau(p)&=1&& (p\in I\text{ a rank-one projection}).
 \end{aligned}
\]
\end{theorem}

\begin{proof}
Let $q$ be multiplication by the indicator of $[1,2]$.
The multiplier $B_t=e^{tP^2}q$ is bounded, so $q=B_tC_t\in I$.
Choose an orthonormal basis of $q\mathcal H$ and let $V$ be its
unilateral shift, extended by zero on $(1-q)\mathcal H$.
Then $V=qVq\in I$, and
\[
 V^*V=q,\qquad VV^*=q-p,
\]
where $p$ projects to the first basis vector. Thus $p\in I$.
Cyclicity gives $\tau(q)=\tau(q-p)$, whereas linearity gives
$\tau(q)-\tau(q-p)=\tau(p)=1$, a contradiction.
\end{proof}

Linearity is essential: the constant function $A\mapsto1$ is cyclic
and takes the value one on every rank-one projection, but is not a
trace. The theorem also does not exclude a weight defined only on a
smaller commutative algebra of multipliers.

That smaller algebra ceases to suffice when all multipliers and the
modular transform act together. Let $\mathcal M$ be their generated
von Neumann algebra. An operator commuting with all multipliers is
itself a multiplier $M_f$. Commutation with the cosine transform
implies
\[
 (f(P)-f(P'))\cos(4\pi PP')=0
 \quad\hbox{for almost every }(P,P').
\]

One obtains this identity by testing the kernels on bounded
intervals. The zero set of the cosine has measure zero, so $f$ is
constant. Hence $\mathcal M'= \mathbb C$ and
$\mathcal M=B(\mathcal H)$.

A faithful normal semifinite trace on this algebra assigns a common
finite positive value $\kappa$ to rank-one projections. Finiteness
follows from semifiniteness on a minimal projection, and equality
from unitary conjugacy. Normality then gives $\tau(q)=\infty$ and
$\tau(C_t)\geq e^{-4t}\tau(q)=\infty$.
Thus a finite multiplier weight cannot extend to such a trace while
retaining the full modular operator algebra.

There is a related tensor obstruction. A vector in
$\overline{\mathcal H}\otimes\mathcal H$ defines a
Hilbert--Schmidt operator. Contracting two such vectors gives a
trace-class operator. It therefore cannot give $C_t$, even for
$t>0$. A two-boundary Hilbert cup for the bare continuous cylinder
requires a different distributional or relative recipient.
The one-boundary preparations of
Proposition~\ref{mht:prop:gravity-prepared-sewing} remain valid.

\section{A positivity test for the regularized vacuum sum}

A different obstruction concerns a particular scalar modular sum.
For $c>1$, take the vacuum character
$q^{-c/24}\prod_{n\geq2}(1-q^n)^{-1}$ and sum its absolute square
over $\Gamma_\infty\backslash\operatorname{PSL}_2(\mathbb Z)$.
Fix the \href{https://arxiv.org/abs/0712.0155}{Maloney--Witten
regularization} by analytic continuation of the
Poincar\'e exponent to $1/2$.

Its large imaginary-modulus expansion is
\[
 Z_c(\tau)=|\chi_{\mathrm{vac}}(\tau)|^2+
 \frac{-6+c_1y^{-1}+c_2y^{-2}+O(y^{-3})}{|\eta(\tau)|^2},
 \qquad \tau=x+iy,
\]
with fixed $x$, where, on putting $u=(c-13)/12$,
\[
 \begin{aligned}
 c_1&=\frac{(\pi^3-6\pi)(c-13)}{9\zeta(3)},\\
 c_2&=\frac{8\pi^2((\pi^4-90)u^2+\pi^4+45)}{135\zeta(5)}>0.
 \end{aligned}
\]

These coefficients belong to the fixed regularized sum. To see the
seed signs, write $\kappa_0=(c-1)/12$. Multiplication of the vacuum
term by $|\eta|^2$ gives
$e^{2\pi\kappa_0y}(1-q)(1-\bar q)$. The four Poincar\'e seeds are
therefore $(\kappa_0,0)$, $(\kappa_0-2,0)$ and
$(\kappa_0-1,\pm1)$, with coefficients $+,+,-,-$.

For a seed $(\kappa,0)$, its continued constant Fourier mode, after
removing the identity term and dividing by $y^{1/2}$, is
\[
 -1+\frac{2\pi^3\kappa}{3\zeta(3)}y^{-1}
      +\frac{4\pi^6\kappa^2}{135\zeta(5)}y^{-2}+O(y^{-3}).
\]
For $(\kappa,\pm1)$ it is
\[
 2+\frac{4\pi\kappa}{\zeta(3)}y^{-1}
      +\frac{4\pi^2(2\kappa^2-1)}{3\zeta(5)}y^{-2}+O(y^{-3}).
\]
These are the constant-mode Fourier expansions of that analytic
continuation. The other Fourier modes are exponentially small at
fixed $x$. Substitution gives $-6,c_1,c_2$ above.
The seed signs and this regularization are part of the question
about a spectral measure.

On the imaginary axis, put $\beta=2\pi y$ and define $d_N$ by
\[
 \prod_{n\geq2}(1-q^n)^{-2}=\sum_{N\geq0}d_Nq^N.
\]
Then
\begin{equation}\label{mht:eq:gravity-vacuum-asymptotic}
 Z_c(\beta)=e^{\beta c/12}\sum_{N\geq0}d_Ne^{-\beta N}
              +e^{\beta/12}\bigl(-6+O(\beta^{-1})\bigr).
\end{equation}
The first coefficients are $1,0,2,2,5,6,13$.

\begin{proposition}\label{mht:prop:gravity-positive-spectrum}
Suppose the function in \eqref{mht:eq:gravity-vacuum-asymptotic} is
the Laplace transform of a positive Borel measure. Assume its energy
support is bounded below and its exponential weighting is finite.
Then $c=12N+1$ for an integer $N\geq5$.
This necessary condition is not a sufficiency statement.
\end{proposition}

\begin{proof}
First suppose a positive measure $\mu$ has
$e^{\beta a}\int e^{-\beta E}\,d\mu(E)\to d<\infty$.
Any mass below $a$ would make this limit infinite. There is thus no
such mass. Dominated convergence after one fixed exponential
weighting shows that the limit is the mass of the atom at $a$.

Apply this observation successively to
\eqref{mht:eq:gravity-vacuum-asymptotic}. It forces actual atoms of
masses $d_N$ at $E_N=N-c/12$ for every
$N<(c-1)/12$. Subtracting these atoms preserves positivity.
At the threshold $E=-1/12$, the remaining rescaled limit is $-6$
unless $m=(c-1)/12$ is an integer. A negative limit is impossible.

If $m$ is an integer, the limit is $d_m-6$, so it must be
nonnegative. The values $d_0,\ldots,d_4$ are all less than six,
while $d_5=6$. Necessarily $m\geq5$, as claimed.
\end{proof}

For example, $c=25$ gives a remaining threshold mass $2-6=-4$.
The cases $c=61$ and $c=73$ give zero and seven, respectively, and
this argument does not rule them out. The test permits continuous
positive measures. It assumes neither a discrete spectrum nor a
previously chosen vacuum subspace. It does not apply unchanged to
signed measures or to a different finite renormalization of the sum.

\section{State normalization and physical reconstruction}

Even a finite algebra can have several positive sewing theories
before its cap is chosen. Let
$A=\mathbb Ce_1\oplus\mathbb Ce_2$, with
$e_ie_j=\delta_{ij}e_i$ and $e_i^*=e_i$. For $0<w<1$, set
\[
 \varepsilon_w(a_1e_1+a_2e_2)=wa_1+(1-w)a_2.
\]

This is a positive normalized trace. Write $p_1=w$, $p_2=1-w$.
The Frobenius pairing has inverse
$\sum_i e_i\otimes e_i/p_i$. Its connected genus-$g$ amplitude is
\[
 Z_g=\sum_{i=1}^2p_i^{\,1-g}.
\]

Indeed, each handle multiplies the $i$th idempotent by $p_i^{-1}$,
and the final cap contributes $p_i$. Thus the sphere and torus
amplitudes are always one and two, whereas genus two distinguishes
$w=1/2$ from $w=1/3$: the values are $4$ and $9/2$.
This is freedom in choosing a cap, not nonuniqueness after the
Frobenius structure has been fixed.

For a gravity comparison, the analogous data must come from the
bulk construction. One needs the gauge quotient with its real
structure, the physical pairing, the cylinder energy and the
renormalized insertion maps. A formal vacuum module or an equality
of oscillator characters does not specify these analytic objects.
Conversely, the positive prepared amplitudes constructed here do
not depend on identifying every bulk state with a boundary vector.

\begin{problem}\label{mht:prob:gravity-physical-comparison}
Construct the comparison from the source-dependent bulk BV/BFV
complex on the solid torus to the boundary state complex. Preserve
the moving boundary condition and both corner variations. Define
the physical gauge quotient and its pairing, and compare its
Euclidean cut with the Lorentzian state space. On this domain,
construct the current, ghost and line insertions with their
renormalized collision products. Determine the coupling and
cylinder-energy normalization from the same theory. Require the
comparison to commute with sewing and with the specified real
structure.
\end{problem}

An obstruction to this problem must survive the admitted local
counterterms and changes of regulator on its physical carrier.
The nonexistence of the ordinary shuffle coproduct and of the finite
continuous-cylinder trace have narrower, explicit domains.
They determine which coefficient or state constructions cannot
serve unchanged in a comparison. The rational Hall tensor and the
positive prepared sewing amplitudes exhibit alternatives whose
physical realization remains a separate question.
